\subsection{অক্টাল থেকে দশমিক}
বিন্দুর বামে ওজন $b^0,b^1,\ldots$; ডানে $b^{-1},b^{-2},\ldots$। প্রতিটি অঙ্কের সঙ্গে তার ওজন গুণ করে অবদানগুলো যোগ করো।
\par\noindent\begin{minipage}{\linewidth}
\convExample{১}{পূর্ণসংখ্যা}
$ (275)_{8} $-কে ১০ ভিত্তিতে রূপান্তর করো।\par
\par\noindent\begin{minipage}{\linewidth}\centering
\phantomsection\label{VC033}
\begin{chThreeDiagram}\begin{tikzpicture}[x=1mm,y=1mm]
\node[vc digit] (d0) at (0,0) {$ 2 $};
\node[vcnum] (w0) at (0,-10) {$ 8^{2} $};
\node[vcnum] (v0) at (0,-25) {$ 128 $};
\draw[vc arrow] (d0.south)--(w0.north);
\draw[vc arrow] (w0.south)--(v0.north);
\node[vc digit] (d1) at (19,0) {$ 7 $};
\node[vcnum] (w1) at (19,-10) {$ 8^{1} $};
\node[vcnum] (v1) at (19,-25) {$ 56 $};
\draw[vc arrow] (d1.south)--(w1.north);
\draw[vc arrow] (w1.south)--(v1.north);
\node[vc digit] (d2) at (38,0) {$ 5 $};
\node[vcnum] (w2) at (38,-10) {$ 8^{0} $};
\node[vcnum] (v2) at (38,-25) {$ 5 $};
\draw[vc arrow] (d2.south)--(w2.north);
\draw[vc arrow] (w2.south)--(v2.north);
\end{tikzpicture}\end{chThreeDiagram}
\end{minipage}\par

\[\begin{aligned}V&=2\times 8^{2}+7\times 8^{1}+5\times 8^{0}\\&=128+56+5=189.\end{aligned}\]
\[\therefore\ (275)_{8}=(189)_{10}.\]
\end{minipage}\par
\par\noindent\begin{minipage}{\linewidth}
\convExample{২}{শুধু ভগ্নাংশ}
$ (0.47)_{8} $-কে ১০ ভিত্তিতে রূপান্তর করো।\par
\par\noindent\begin{minipage}{\linewidth}\centering
\phantomsection\label{VC034}
\begin{chThreeDiagram}\begin{tikzpicture}[x=1mm,y=1mm]
\node[vc digit] (d0) at (0,0) {$ 0 $};
\node[vcnum] (w0) at (0,-10) {$ 8^{0} $};
\node[vcnum] (v0) at (0,-25) {$ 0 $};
\draw[vc arrow] (d0.south)--(w0.north);
\draw[vc arrow] (w0.south)--(v0.north);
\node[vc digit] (d1) at (27,0) {$ 4 $};
\node[vcnum] (w1) at (27,-10) {$ 8^{-1} $};
\node[vcnum] (v1) at (27,-25) {$ \frac{1}{2} $};
\draw[vc arrow] (d1.south)--(w1.north);
\draw[vc arrow] (w1.south)--(v1.north);
\node[vc digit] (d2) at (46,0) {$ 7 $};
\node[vcnum] (w2) at (46,-10) {$ 8^{-2} $};
\node[vcnum] (v2) at (46,-25) {$ \frac{7}{64} $};
\draw[vc arrow] (d2.south)--(w2.north);
\draw[vc arrow] (w2.south)--(v2.north);
\node[vc digit] at (13.5,0) {$.$};
\end{tikzpicture}\end{chThreeDiagram}
\end{minipage}\par

\[\begin{aligned}V&=0\times 8^{0}+4\times 8^{-1}+7\times 8^{-2}\\&=\frac{1}{2}+\frac{7}{64}=0.609375.\end{aligned}\]
\[\therefore\ (0.47)_{8}=(0.609375)_{10}.\]
\end{minipage}\par
\par\noindent\begin{minipage}{\linewidth}
\convExample{৩}{মিশ্র সংখ্যা}
$ (356.24)_{8} $-কে ১০ ভিত্তিতে রূপান্তর করো।\par
\par\noindent\begin{minipage}{\linewidth}\centering
\phantomsection\label{VC035}
\begin{chThreeDiagram}\begin{tikzpicture}[x=1mm,y=1mm]
\node[vc digit] (d0) at (0,0) {$ 3 $};
\node[vcnum] (w0) at (0,-10) {$ 8^{2} $};
\node[vcnum] (v0) at (0,-25) {$ 192 $};
\draw[vc arrow] (d0.south)--(w0.north);
\draw[vc arrow] (w0.south)--(v0.north);
\node[vc digit] (d1) at (19,0) {$ 5 $};
\node[vcnum] (w1) at (19,-10) {$ 8^{1} $};
\node[vcnum] (v1) at (19,-25) {$ 40 $};
\draw[vc arrow] (d1.south)--(w1.north);
\draw[vc arrow] (w1.south)--(v1.north);
\node[vc digit] (d2) at (38,0) {$ 6 $};
\node[vcnum] (w2) at (38,-10) {$ 8^{0} $};
\node[vcnum] (v2) at (38,-25) {$ 6 $};
\draw[vc arrow] (d2.south)--(w2.north);
\draw[vc arrow] (w2.south)--(v2.north);
\node[vc digit] (d3) at (65,0) {$ 2 $};
\node[vcnum] (w3) at (65,-10) {$ 8^{-1} $};
\node[vcnum] (v3) at (65,-25) {$ \frac{1}{4} $};
\draw[vc arrow] (d3.south)--(w3.north);
\draw[vc arrow] (w3.south)--(v3.north);
\node[vc digit] (d4) at (84,0) {$ 4 $};
\node[vcnum] (w4) at (84,-10) {$ 8^{-2} $};
\node[vcnum] (v4) at (84,-25) {$ \frac{1}{16} $};
\draw[vc arrow] (d4.south)--(w4.north);
\draw[vc arrow] (w4.south)--(v4.north);
\node[vc digit] at (51.5,0) {$.$};
\end{tikzpicture}\end{chThreeDiagram}
\end{minipage}\par

\[\begin{aligned}V&=3\times 8^{2}+5\times 8^{1}+6\times 8^{0}+2\times 8^{-1}+4\times 8^{-2}\\&=192+40+6+\frac{1}{4}+\frac{1}{16}=238.3125.\end{aligned}\]
\[\therefore\ (356.24)_{8}=(238.3125)_{10}.\]
\end{minipage}\par
